\section{Order of magnitude per PDS readout channel}
\label{sec:perchannel}

The requested comparison is the time-averaged \textbf{primary scintillation
PE/s per electronic readout channel}. The cited designs have
\textbf{32 channels in ProtoDUNE-VD}, \textbf{1344 in one FD-VD module},
\textbf{8400 in full ND-LAr}, and \textbf{384 in ND 2$\times$2}
\cite{pd_channels,vd_tdr,nd_cdr,nd_channels,nd_operation}.
Full ND-LAr and the 2$\times$2 prototype are separate detectors.

\begin{center}
\small
\begin{tabular}{lrrr}
\toprule
Readout group & Channels & $^{39}$Ar & Cosmic tracks \\
 & in group & PE/s/channel & PE/s/channel \\
\midrule
ProtoDUNE-VD cathode XA & 16 & $7.7\times10^4$ & $1.2\times10^7$ \\
ProtoDUNE-VD wall XA & 16 & $3.9\times10^4$ & $6.1\times10^6$ \\
\addlinespace
FD-VD cathode XA & 640 & $2.2\times10^5$ & $8.5$ \\
FD-VD long-wall XA & 640 & $1.1\times10^5$ & $4.3$ \\
FD-VD short-wall XA & 64 & $1.1\times10^5$ & $4.2$ \\
\addlinespace
Full ND-LAr ArCLight & 4200 & $44$ & $90$ \\
Full ND-LAr LCM & 4200 & $133$ & $271$ \\
\addlinespace
ND 2$\times$2 ArCLight & 192 & $15$ & $12$ \\
ND 2$\times$2 LCM & 192 & $48$ & $39$ \\
\bottomrule
\end{tabular}
\end{center}
These are \textbf{conditional reference estimates, not measured rates}.
VD rows use representative XA positions, with equal sharing between two
channels per XA; they are not averages over the installed arrays. ND rows
are technology averages under a uniform first-hit wall-flux approximation.
Cosmic rows use the site and track assumptions in Section~\ref{sec:cosmicnumbers},
with no beam light. Dark avalanches, correlated SiPM noise, and other
radioactivity are excluded.

For ND, the means over \emph{all} light channels are
\begin{equation*}
\begin{array}{lcc}
 & {}^{39}\mathrm{Ar} & \text{cosmic-track reference}\\
\text{Full ND-LAr:} & \boxed{89\ \mathrm{PE/s/channel}} & 180\ \mathrm{PE/s/channel}\\[2pt]
\text{ND 2$\times$2:} & \boxed{31\ \mathrm{PE/s/channel}} & 26\ \mathrm{PE/s/channel}.
\end{array}
\end{equation*}
The radioactive-light scales are therefore $10^4$--$10^5$ PE/s/channel for
ProtoDUNE-VD, $10^5$ for FD-VD, $10^2$ for full ND-LAr, and a few $10^1$
for 2$\times$2. Optical segmentation, source mass per compartment, aperture,
and PDE determine the light per channel.

\subsection{Numerical inputs and channel sharing}

Use atmospheric-argon activity $S_{39}=\SI{0.964}{Bq/kg}$,
$\rho=\SI{1395}{kg/m^3}$, and the explicit reference yield
\begin{equation}
\overline N_{\gamma,39}^{\mathrm{ref}}
=(\SI{0.220}{MeV})(20{,}000\ \mathrm{photons/MeV})=4400\ \mathrm{photons/decay}.
\label{eq:yieldref}
\end{equation}
The mean beta energy is quoted in the DUNE interim design report
\cite{dune_idr}; the VD TDR uses a total yield of 20,000 photons/MeV for
its Xe-doped simulation~\cite{vd_tdr}. Their product is a reference
normalization, not a measured beta-spectrum-averaged yield in each detector.
Scale each $^{39}$Ar entry by its actual $\overline N_{\gamma,39}/4400$.

For VD, adopt a whole-device XA PDE of $\epsilon_{\mathrm{XA}}=0.045$
\cite{fdvd_pde} and calculate
\begin{equation}
\boxed{R_{39,\mathrm{ch}}=a_{39}G_{\mathrm{XA}}
\overline N_{\gamma,39}\epsilon_{\mathrm{XA}}/2.}
\label{eq:VDchannel}
\end{equation}
The factor $1/2$ shares detected light across two readout channels; it
leaves the entrance area and emitting drift regions unchanged.
Each row scales by $\epsilon_{\mathrm{XA}}/0.045$ and by changes in optical
acceptance. Direct propagation here is unobstructed and lossless.

Cathode entries use the centered double-sided geometry derived below.
The wall benchmark is single-sided and sees one drift region: for
ProtoDUNE use normal depth 3.0 m, tangential dimensions $6.8\times3.5$ m,
and a centered aperture. For FD long/short walls use normal depths
13.5/60 m, tangential dimensions $60\times6.5$/$13.5\times6.5$ m, and
height 4.5 m above the cathode (1.25 m above the half-drift center).
This assumes optical separation at the cathode. The resulting
$G_{\mathrm{XA}}$ values are 0.290, 0.862 and 0.842 m$^3$, respectively.
Actual wall and cathode placements require their own integrals.

For ND, adopt whole-panel PDEs 0.002 (ArCLight) and 0.006 (LCM) from the
prototype design summary~\cite{nd_status}, explicitly as reference values
for both sizes. Divide the light assigned to each technology by its own
channel count. The coverage, compartment masses, and channel derivation
are given in Section~\ref{sec:nd2x2}.

\subsection{If ``activity per channel'' means Bq}

Dividing total modeled $^{39}$Ar decays by the readout count gives the
following bookkeeping normalization:
\begin{center}
\small
\begin{tabular}{lrrr}
\toprule
Detector & Source mass (t) & Total decays/s & Decays/s divided by channels \\
\midrule
ProtoDUNE-VD & 199 & $1.92\times10^5$ & $6.00\times10^3$ \\
FD-VD & 14,689 & $1.42\times10^7$ & $1.05\times10^4$ \\
Full ND-LAr & 146.5 & $1.41\times10^5$ & $16.8$ \\
ND 2$\times$2 & 2.4 & $2.31\times10^3$ & $6.03$ \\
\bottomrule
\end{tabular}
\end{center}
VD masses follow the reference source boxes; ND-LAr follows 35 active
$1\times1\times3$ m modules~\cite{nd_cdr}. These are not cryostat totals
or fiducial selection masses. The quotient is \emph{not} a channel event
rate: a decay can illuminate several channels or produce no detected PE.
The optical calculation above supplies the detected-light comparison.

\clearpage
